% Manuscript-ready Theory section fragment.
% Requires amsmath and amssymb. Intended for direct inclusion via \input{}.

\section{Theory: consistency gaps as composite finite-training interventions}
\label{sec:theory}

\subsection{Ideal target and endpoint-invariance intuition}
\label{sec:theory-ideal-endpoint}

Let $x_t$ denote a state on a coupled trajectory and let
$\mathcal{T}_{t\rightarrow r}$ transport that state from $t$ to $r\leq t$.
We call a gap schedule $S$ \emph{legal} on a support $\mathcal{X}$ if every
supported time $t$ is connected to a common anchor $t_\star$ by a finite chain
$t=t_0>t_1>\cdots>t_m=t_\star$ of scheduled edges. Suppose that (i) the
transport is exact and compositional,
$\mathcal{T}_{r\rightarrow s}\!\circ\mathcal{T}_{t\rightarrow r}
=\mathcal{T}_{t\rightarrow s}$; (ii) all scheduled edges are covered by the
population objective; (iii) the function class can attain zero consistency
residual on those edges; (iv) optimization reaches such a population solution;
and (v) the mapping is fixed at the common anchor by
$F(x,t_\star)=b(x)$. Then zero residual along the schedule telescopes to
\begin{equation}
  F^\star(x_t,t)
  =F^\star\!\left(\mathcal{T}_{t\rightarrow t_1}(x_t),t_1\right)
  =\cdots
  =b\!\left(\mathcal{T}_{t\rightarrow t_\star}(x_t)\right).
  \label{eq:ideal-endpoint-invariance}
\end{equation}
Consequently, any two legal schedules that share the same transport, support,
and anchor identify the same ideal consistency mapping on their reachable
support. The gap schedule can therefore change the constraints presented during
training without changing the ideal endpoint mapping under these assumptions.
This is an ideal identification statement, not a claim about finite-capacity,
finite-data, or incompletely optimized models; without a common anchor,
connectivity, support coverage, or an attainable zero-residual solution,
schedule invariance need not hold.

\subsection{Exact ECT composite intervention}
\label{sec:theory-composite-intervention}

The preceding ideal invariance does not make the ECT training intervention
scalar. For one coupled ECT sample, let
$x_t=y+t\varepsilon$, $x_r=y+r\varepsilon$, $\Delta=t-r>0$, and define
\begin{equation}
  D_t=f_\theta(x_t,t),\qquad
  T_r=\operatorname{sg}\!\left[f_\theta(x_r,r)\right],\qquad
  e_r=D_t-T_r,qquad
  \phi(e)=\lVert e\rVert_2 .
  \label{eq:ect-residual}
\end{equation}
Here $\operatorname{sg}$ preserves the target's forward value but suppresses its
parameter derivative. Formally, for a two-parameter lifted loss
$\widetilde\ell(\theta,\bar\theta)$, the one-sided stop-gradient training
derivative is
\begin{equation}
  \mathcal{D}^{\mathrm{sg}}_\theta\ell
  :=\left.
  \partial_\theta\widetilde\ell(\theta,\bar\theta)
  \right|_{\bar\theta=\theta},
  \label{eq:sg-training-derivative}
\end{equation}
which differs from the ordinary derivative along the parameter diagonal. The
ECT objective and its one-sided training derivative are, for $e_r\neq0$,
\begin{equation}
  \ell(r,\Delta)=\frac{\phi(e_r)}{\Delta},
  \qquad
  g(r,\Delta):=\mathcal{D}^{\mathrm{sg}}_\theta\ell(r,\Delta)
  =\frac{1}{\Delta}J_t^\top u_r,
  \qquad
  J_t:=\partial_\theta D_t,
  \quad
  u_r:=\frac{e_r}{\lVert e_r\rVert_2}.
  \label{eq:exact-ect-gradient}
\end{equation}
Consider a baseline gap $\Delta_1$ with target
$r_1=t-\Delta_1$ and a probe gap $\Delta_g$ with target
$r_g=t-\Delta_g$, evaluated at the same parameter value and with the same
online realization. Setting $s:=\Delta_1/\Delta_g$, direct substitution gives
the exact finite-gap decompositions
\begin{align}
  \ell_g
  &=s\ell_1
    +\underbrace{\frac{\phi(e_{r_g})-\phi(e_{r_1})}{\Delta_g}}
      _{\tau^{(\ell)}_{\mathrm{tar}}},
  \label{eq:exact-ect-loss-decomposition}\\
  g_g
  &=s g_1
    +\underbrace{\frac{1}{\Delta_g}J_t^\top
      \left(u_{r_g}-u_{r_1}\right)}
      _{\tau^{(g)}_{\mathrm{tar}}}.
  \label{eq:exact-ect-gradient-decomposition}
\end{align}
We call $s\ell_1$ and $s g_1$ the \emph{denominator-rescaling component}, and
$\tau^{(\ell)}_{\mathrm{tar}}$ and $\tau^{(g)}_{\mathrm{tar}}$ the
\emph{target-endpoint component}. Their joint change is the \emph{ECT composite
gap intervention}. Equations~\eqref{eq:exact-ect-loss-decomposition}--
\eqref{eq:exact-ect-gradient-decomposition} are algebraic identities: they
require neither a small-gap approximation nor an optimizer model. In
particular, only the denominator-rescaling component is exactly scalar; the
total training derivative is scalar-equivalent to the baseline only when its
target-endpoint component vanishes or happens to align with the baseline
derivative.

\subsection{Exact factorial identities}
\label{sec:theory-factorial-identities}

The composite intervention yields an exact objective-gradient factorialization.
For a matched sample, define the four cells
\begin{equation}
  A=(r_1,\Delta_1),\qquad
  B=(r_g,\Delta_g),\qquad
  C=(r_g,\Delta_1),\qquad
  D=(r_1,\Delta_g),
  \label{eq:factorial-cells}
\end{equation}
where the first coordinate selects the detached target endpoint and the second
selects the explicit denominator. Because the denominator multiplies the
objective and its one-sided training derivative externally, the cell values
satisfy
\begin{align}
  \ell_B&=s\ell_C,&
  \ell_D&=s\ell_A,&
  \ell_B-\ell_D&=s(\ell_C-\ell_A),
  \label{eq:factorial-loss-identities}\\
  g_B&=s g_C,&
  g_D&=s g_A,&
  g_B-g_D&=s(g_C-g_A).
  \label{eq:factorial-gradient-identities}
\end{align}
The corresponding factorial interactions are therefore not free residuals but
obey
\begin{equation}
  I_\ell:=\ell_B-\ell_C-\ell_D+\ell_A
  =(s-1)(\ell_C-\ell_A),
  \qquad
  I_g:=g_B-g_C-g_D+g_A
  =(s-1)(g_C-g_A).
  \label{eq:exact-factorial-interactions}
\end{equation}
For an unclipped global gap multiplier, $s=1/g$ is common to all samples and
these identities also hold after minibatch aggregation; with clipping they
remain exact per sample but generally do not reduce to one batch-wide scalar.
These identities factorize the intervention at the objective-gradient level;
they do not imply an additive or seed-invariant factorization of finite-training
quality. Once the four arms follow separate parameter trajectories, the matched
quantities in Eqs.~\eqref{eq:factorial-loss-identities}--
\eqref{eq:exact-factorial-interactions} are no longer held fixed, so the theory
motivates the factorial experiment without predetermining its finite-training
quality contrasts.
